\section{Background and Motivation}
Hospital Management Systems (HMS) store and process sensitive information such as medical records, lab results, diagnoses, prescriptions, and personal patient data. In a hospital, authorization is not only about security. It also affects clinical response time. Patient data must be protected because it is confidential, but doctors and nurses must be able to access the right information immediately, especially in emergencies. This makes access control in hospitals a constant trade-off between confidentiality and availability.
\vspace{0.5cm}


Many hospitals still rely on Role-Based Access Control (RBAC) because it fits the hospital structure and is easy to manage. Roles such as \textit{Doctor}, \textit{Nurse}, \textit{Lab Technician}, and \textit{Pharmacist} map naturally to daily responsibilities. However, traditional RBAC is often too rigid for real hospital workflows. A doctor should not automatically have access to every patient record just because the role is ``Doctor''. Access usually depends on the ward, the care-team, and the current assignment. Similarly, a nurse may update patient observations during an active shift, but the same action should be restricted outside duty hours. These are common dynamic situations, and RBAC alone does not capture them well without adding many special-case roles or manual exceptions.
\vspace{0.5cm}


Attribute-Based Access Control (ABAC) can describe these situations better because it uses context such as ward membership, shift status, record sensitivity, and emergency conditions. For example, in an emergency case, a doctor on call may need immediate access even if the patient is not yet assigned to that care-team. ABAC supports this type of decision logic. However, a pure ABAC approach is difficult to operate in practice. Policies can grow quickly, and the system becomes harder to manage, test, and audit. In a long-running hospital system, this operational complexity is a real problem.
\vspace{0.5cm}


For these reasons, this thesis implements a Hybrid RBAC/ABAC model in a role-centric manner for a microservices-based HMS. RBAC is used as the stable baseline because it matches the hospital hierarchy and is familiar to administrators. ABAC is used as a constraint layer to refine RBAC decisions using dynamic context, such as ward, care-team, shift, emergency flags, and record sensitivity. This role-centric hybrid approach keeps policies readable and maintainable while still supporting the flexibility required in hospital operations.
\vspace{0.5cm}


This design also matches microservices deployment. In a microservices HMS, data and functions are distributed across services, for example medical records, laboratory results, prescriptions, and billing. If each service enforces authorization differently, inconsistencies and security gaps are likely to appear. A hybrid approach with a clear decision flow helps enforce consistent access rules across services while keeping the model practical for maintenance.

\section{Problem Statement}
A microservices-based HMS needs an authorization mechanism that protects confidential patient data, supports immediate clinical access when needed, and remains manageable as the system evolves. Traditional RBAC is simple and widely used, but it cannot handle dynamic hospital context well without becoming overly permissive or overly complex. Pure ABAC is flexible, but it can be difficult to maintain and audit due to policy growth and increasing rule complexity.
\vspace{0.5cm}


In addition, the HMS environment changes over time. New services may be added, workflows may be updated, and policies may increase in number. Therefore, the system should maintain consistent enforcement across services, keep decision latency acceptable under policy growth, and provide a way to detect and handle policy conflicts to avoid contradictory authorization outcomes.

\subsection{Research Objectives}
This thesis aims to (i) define hospital roles, resources, actions, and key access constraints; (ii) build an RBAC baseline for core HMS functions; (iii) refine authorization using role-centric ABAC constraints for dynamic context; (iv) design and implement an enforcement architecture for microservices (e.g., PEP--PDP--PIP); and (v) evaluate correctness, performance, scalability under increasing policy counts, and policy conflict detection.

\section{Research Scope}
This thesis focuses on application-level authorization for an HMS. The scope includes a role-centric Hybrid RBAC/ABAC model where RBAC provides the baseline and ABAC refines decisions using context. The thesis covers typical HMS operations such as reading and updating medical records, accessing lab results, and prescription-related actions. Infrastructure-level security topics such as advanced network hardening and low-level storage encryption are not the main focus, except when they directly support the access control discussion.

\section{Methodology and Evaluation Plan}

\subsection{Methodology}
This thesis follows a practical development process for building and validating a role-centric Hybrid RBAC/ABAC solution in a microservices-based HMS. First, hospital workflows are analyzed to identify the main users, resources, and actions, together with realistic access constraints (e.g., ward restriction, care-team membership, duty shift, emergency access, and record sensitivity). Based on these requirements, an RBAC baseline is defined to reflect the hospital’s natural hierarchy and job functions. After that, ABAC constraints are designed as a refinement layer on top of RBAC to capture dynamic context without making the overall policy set difficult to operate.
\vspace{0.5cm}


Next, the thesis designs an authorization architecture suitable for microservices, following the PEP--PDP--PIP pattern. Policies are stored in a central policy store, enforcement is placed at service boundaries, and attributes are retrieved from relevant sources to form complete request context. Finally, the model is implemented and integrated into selected services to support consistent enforcement and auditing.

\subsection{Evaluation Plan}
The evaluation focuses on whether the implemented model is correct, practical, and scalable under realistic policy growth. The thesis evaluates correctness using scenario-based test cases with expected decisions. It evaluates system performance by measuring decision latency and throughput under workload. It studies scalability by increasing the number of policies and observing how evaluation time and resource usage change. In addition, the thesis evaluates policy conflict detection to ensure policy evolution does not create contradictory outcomes.

\begin{table}[H]
\centering
\renewcommand{\arraystretch}{1.2}
\setlength{\tabcolsep}{6pt}
\begin{tabularx}{0.98\textwidth}{|
>{\centering\arraybackslash}p{0.9cm}|
>{\raggedright\arraybackslash}p{3.2cm}|
>{\raggedright\arraybackslash}X|
>{\raggedright\arraybackslash}p{4.2cm}|
}
\hline
\textbf{ID} & \textbf{Evaluation aspect} & \textbf{What is evaluated} & \textbf{How it is measured} \\
\hline
1 & Correctness &
Whether authorization decisions match the intended policies across realistic HMS scenarios. &
Scenario-based test matrix (expected vs. actual permit/deny); report false-permit and false-deny cases. \\
\hline
2 & Performance &
The runtime cost of authorization evaluation in microservices (PDP decision time and end-to-end enforcement time). &
Decision latency (ms) and throughput (requests/sec) under controlled workloads; compare average and tail latency (e.g., P95/P99) if available. \\
\hline
3 & Scalability (policy growth) &
How decision time and resource usage change when the number of policies increases. &
Run the same workload with increasing policy counts (e.g., 100 / 500 / 1000); measure latency, throughput, and memory usage. \\
\hline
4 & Scalability (context size) &
Sensitivity to the number of attributes and context elements used in evaluation. &
Increase the number of attributes in the request context (subject/resource/environment) and measure the effect on decision latency and memory usage. \\
\hline
5 & Policy conflict detection &
Ability to detect and handle contradictory or overlapping policies during policy updates. &
Number of conflicts found and detection time under different policy set sizes; describe conflict types and resolution strategy (e.g., deny-overrides, specificity). \\
\hline

\end{tabularx}
\end{table}

\begin{table}[H]
    \centering
\renewcommand{\arraystretch}{1.2}
\setlength{\tabcolsep}{6pt}
\label{tab:evaluation-plan}
\begin{tabularx}{0.98\textwidth}{|
>{\centering\arraybackslash}p{0.9cm}|
>{\raggedright\arraybackslash}p{3.2cm}|
>{\raggedright\arraybackslash}X|
>{\raggedright\arraybackslash}p{4.2cm}|
}
\hline
\textbf{ID} & \textbf{Evaluation aspect} & \textbf{What is evaluated} & \textbf{How it is measured} \\
\hline
6 & Consistency across services &
Whether enforcement remains consistent across multiple microservices that share policies. &
Cross-service test cases using the same user/context; check decision consistency and audit logs for the same policy intent. \\
\hline
\end{tabularx}
    \caption{Evaluation plan for the role-centric Hybrid RBAC/ABAC model}
\label{tab:evaluation-plan}
\end{table}

\section{Thesis Structure}
The thesis is organized as follows:
\begin{itemize}
    \item \textbf{Chapter 2:} presents related background: authentication and authorization, RBAC and ABAC models, hybrid RBAC/ABAC, and microservices concepts related to access control.
    \item \textbf{Chapter 3:} describes system analysis and design: requirements, use cases, the proposed hybrid model design, system architecture, and data/policy modeling.
    \item \textbf{Chapter 4:} presents implementation and evaluation: how the services and authorization logic are implemented, how attributes are collected, the evaluation setup, and the results with discussion.
    \item \textbf{Chapter 5:} concludes the thesis and suggests future improvements.
\end{itemize}

\input{chapters/RelatedWorks}
